% =======================================================================================
% MODELO DE TCC - INSTITUTO FEDERAL DE SERGIPE (IFS)
% =======================================================================================
% Autores do Modelo (discentes de Ciência da Computação - IFS, Campus Itabaiana):
%   Mikael dos Santos Gois
%   Talyson Rodrigues Silva Nascimento
%   Kauan Gutierrez Santos Estácio
% Modelo Base: abnTeX2, adaptado para as normativas internas de TCC do IFS (Art. 31 e 34)
% Créditos dos modelos base (abnTeX2, DECOM/UFOP e Diogo D. Moreira) e das imagens: ver NOTICE
% Licença: CC BY-SA 4.0 (ver LICENSE)
% Este arquivo faz parte de um modelo opensource para estudantes do IFS.
% Não remova os créditos dos autores originais e do modelo base.
% Este modelo é fornecido "no estado em que se encontra", sem garantias de qualquer tipo.
% =======================================================================================

% [O que aparece no PDF Bookmark] {O que aparece impresso no documento}

\chapter[Conclusão]{CONCLUSÃO}
\label{sec:conclusao}

% EXEMPLO DIDÁTICO DE CONCLUSÃO
% Todos os resultados mencionados são fictícios.
% Substitua este conteúdo pelos dados reais do TCC.

Esta conclusão integra um exemplo didático, elaborado com dados fictícios, para demonstrar como sintetizar os resultados e explicitar seu alcance em um trabalho acadêmico. Na versão definitiva, esta seção deve responder ao problema de pesquisa e retomar o objetivo
geral apresentado na Introdução.

No exemplo desenvolvido nos capítulos anteriores, o propósito da avaliação foi comparar duas
configurações de uma aplicação computacional quanto ao tempo médio de execução e ao consumo médio
de memória. Os valores hipotéticos apresentados indicaram uma redução de 25\% no tempo de execução
e de aproximadamente 12,9\% no consumo de memória na configuração com otimização. Essas diferenças servem apenas para ilustrar a síntese de resultados; não representam evidências experimentais reais.

A interpretação do exemplo mostra como uma comparação quantitativa pode subsidiar a avaliação
de uma solução computacional. Entretanto, a identificação de diferenças entre duas medições,
isoladamente, não demonstra que uma alternativa será superior em qualquer ambiente ou carga de
trabalho. A conclusão de um estudo real deve permanecer limitada aos cenários, às métricas e
aos procedimentos efetivamente avaliados.

Entre as limitações do exemplo estão a ausência de informações sobre o número de repetições,
a variabilidade das medições e o controle das condições de execução. Por esse motivo, não é
possível avaliar a consistência estatística das diferenças apresentadas nem generalizar o
comportamento observado.

Como possibilidade de aprofundamento em uma pesquisa real, recomenda-se ampliar os cenários
de teste, executar medições repetidas, registrar as características do ambiente e analisar a
dispersão dos resultados. Essas ações podem contribuir para uma avaliação mais robusta
da solução estudada.

Em síntese, este exemplo demonstra uma estrutura possível para a conclusão: retomada do objetivo,
síntese dos achados, interpretação do seu alcance, reconhecimento das limitações e indicação de
desdobramentos. Ao utilizar o modelo em um TCC, todo o conteúdo ilustrativo deve ser substituído
por conclusões fundamentadas nos dados e nas análises da própria pesquisa.


% -------------------------------------------------------
% DICA PARA O ESTUDANTE — CONCLUSÃO
% -------------------------------------------------------
% A conclusão não deve simplesmente repetir os
% resultados. Sua função é apresentar o significado
% dos achados em relação ao problema e aos objetivos
% definidos na Introdução.
%
% Ao adaptar o exemplo, o estudante deve substituir
% os números fictícios, explicar se os objetivos
% foram atingidos com base nas evidências, identificar
% as contribuições concretas, reconhecer limitações
% e propor possíveis continuações.
